\chapter{Sprint 2 : Module de Découverte et de Mise en Relation}
\minitoc
\newpage

\section{Introduction}
La valeur centrale de GOSPORT réside dans la capacité à connecter des athlètes avec des professionnels du sport (coachs, nutritionnistes) de manière fluide et sécurisée. Ce deuxième sprint implémente le cycle complet de mise en relation : la découverte de profils, la consultation des fiches détaillées, et le flux de demandes de coaching incluant leur acceptation ou leur refus.

À l'issue de ce sprint, un athlète peut parcourir l'annuaire des professionnels disponibles sur la plateforme, consulter leurs profils, et envoyer une demande de suivi. Le professionnel peut alors consulter les demandes en attente et décider de les accepter ou de les refuser.

\section{Backlog du Sprint 2}

\begin{table}[H]
\centering
\renewcommand{\arraystretch}{1.4}
\resizebox{\textwidth}{!}{%
\begin{tabular}{|c|c|p{5cm}|p{6.5cm}|c|}
\hline
\rowcolor[HTML]{DAE8FC}
\textbf{ID} & \textbf{Acteur} & \textbf{User Story} & \textbf{Tâche Technique} & \textbf{Priorité} \\ \hline
2.1 & Athlète & Découvrir et filtrer les coachs disponibles & Filtres par spécialité, tarif et note & Must \\ \hline
2.2 & Tous & Consulter le profil public d'un professionnel & Page de détail du profil professionnel & Must \\ \hline
2.3 & Athlète & Envoyer une demande de coaching & Entité CoachingRequest et workflow de demande & Must \\ \hline
2.4 & Coach & Accepter ou refuser une demande & Mise à jour du statut de la demande & Must \\ \hline
2.5 & Athlète & Envoyer une demande à un nutritionniste & Workflow de demande nutritionnelle & Should \\ \hline
2.6 & Nutritionniste & Répondre à une demande de suivi nutrition & Mise à jour du statut nutritionnel & Should \\ \hline
2.7 & Coach & Dissocier un athlète de son suivi & Action de dissociation sécurisée & Could \\ \hline
\end{tabular}%
}
\caption{Backlog du Sprint 2 -- Module de mise en relation}
\end{table}

\section{Conception et Modélisation}

\subsection{Diagramme de Cas d'Utilisation}
\begin{figure}[H]
\centering
\includegraphics[width=0.88\textwidth]{Sprint2_Cas_Utilisation.png}
\caption{Diagramme de cas d'utilisation -- Sprint 2}
\end{figure}

\subsection{Diagramme de Classes -- Sprint 2}
\begin{figure}[H]
\centering
\includegraphics[width=0.9\textwidth]{Sprint2_Classes.png}
\caption{Diagramme de classes -- Mise en relation}
\end{figure}

Les entités clés de ce sprint sont :\texttt{CoachingRequest} (statut: PENDING \textbar{} ACCEPTED \textbar{} REJECTED) et \texttt{NutritionConnection}. Ces deux entités établissent un lien entre un \texttt{User} (athlète) et un professionnel (coach ou nutritionniste). Le cycle de vie de cet objet -- de sa création à son acceptation -- est central pour ouvrir l'accès aux fonctionnalités suivantes (programmes, nutrition, messagerie).

\subsection{Diagramme de Séquence : Demande de Coaching}
\begin{figure}[H]
\centering
\includegraphics[width=0.88\textwidth]{Sprint2_Sequence_CoachingRequest.png}
\caption{Diagramme de séquence -- Cycle de vie d'une demande de coaching}
\end{figure}

\section{Architecture Technique}


\subsection{Modèle de Données : CoachingRequest}
L'entité \texttt{CoachingRequest} est le pivot de ce sprint. Elle porte un statut énuméré qui suit un cycle de vie strict :

\begin{enumerate}
    \item \textbf{PENDING :} État initial créé par la requête \texttt{POST} de l'athlète.
    \item \textbf{ACCEPTED :} Le coach accepte via \texttt{PATCH} -- la relation de coaching est désormais active.
    \item \textbf{REJECTED :} Le coach refuse -- la demande est archivée sans suppression pour traçabilité.
\end{enumerate}

\subsection{Pagination et Filtrage des Profils}
Le module de recherche des coachs supporte la pagination et plusieurs critères de filtrage afin d'offrir une expérience de recherche pertinente :
\begin{itemize}
    \item \textbf{?specialty=nutrition} : Filtrage par spécialité professionnelle.
    \item \textbf{?maxPrice=100} : Filtrage par tarif mensuel maximum.
    \item \textbf{?page=1\&limit=12} : Pagination des résultats (12 profils par page).
\end{itemize}

\section{Réalisation}

\subsection{Contrôleur de Requêtes de Coaching}
\begin{lstlisting}[language=JavaScript,
    caption=Extrait de CoachingRequestController.ts]
export const createRequest = async (req: Request, res: Response) => {
  const { coachId, message } = req.body;
  const athleteId = req.user.id;

  // Verification : eviter les doublons actifs
  const existing = await requestRepo.findOne({
    where: { athleteId, coachId,
             status: RequestStatus.PENDING }
  });
  if (existing)
    return res.status(409).json({ message: 'Demande deja envoyee' });

  const request = requestRepo.create({
    athleteId, coachId, message,
    status: RequestStatus.PENDING
  });
  await requestRepo.save(request);

  // Notification au coach
  await notifService.notify(coachId,
    'Nouvelle demande de coaching recue');

  return res.status(201).json(request);
};
\end{lstlisting}

\subsection{Composant Discovery -- Angular}
La page de découverte Angular (\texttt{DiscoveryComponent}) charge dynamiquement les profils depuis l'API et applique les filtres côté client ou serveur selon la configuration choisie. Les profils sont présentés sous forme de cartes (cards) interactives affichant la photo de profil, la note moyenne, la spécialité et le tarif mensuel.

\begin{lstlisting}[language=JavaScript, caption=Extrait discovery.component.ts]
export class DiscoveryComponent implements OnInit {
  coaches: Coach[] = [];
  filters = { specialty: '', maxPrice: null };

  ngOnInit() { this.loadCoaches(); }

  loadCoaches() {
    this.coachService.getCoaches(this.filters)
      .subscribe(data => this.coaches = data);
  }

  sendRequest(coachId: number) {
    this.coachingService.sendRequest(coachId, '')
      .subscribe({
        next: () => this.toast.success('Demande envoyee !'),
        error: (e) => this.toast.error(e.error.message)
      });
  }
}
\end{lstlisting}

\section{Tests et Validation}

\begin{table}[H]
\centering
\renewcommand{\arraystretch}{1.3}
\begin{tabular}{|p{5.5cm}|l|l|}
\hline
\rowcolor[HTML]{FFF2CC}
\textbf{Scénario de test} & \textbf{Résultat attendu} & \textbf{Statut} \\ \hline
Lister les coachs (filtres vides) & 200 + liste paginée & \checkmark Conforme \\ \hline
Filtrer par spécialité & 200 + liste filtrée & \checkmark Conforme \\ \hline
Envoyer une demande valide & 201 Created & \checkmark Conforme \\ \hline
Envoi de demande en doublon & 409 Conflict & \checkmark Conforme \\ \hline
Coach accepte une demande & 200 + status ACCEPTED & \checkmark Conforme \\ \hline
Coach refuse une demande & 200 + status REJECTED & \checkmark Conforme \\ \hline
\end{tabular}
\caption{Résultats des tests d'intégration -- Sprint 2}
\end{table}

\section{Bilan du Sprint 2}
Le module de mise en relation est complet. Les athlètes peuvent désormais explorer la place de marché des professionnels, consulter les profils publics et initier des demandes de coaching ou nutritionnelles. Les coachs disposent d'un tableau de bord pour gérer ces demandes entrantes. L'établissement de ces liens est la clé qui déverrouille l'accès aux fonctionnalités des sprints suivants (partage de programmes, nutrition, messagerie).
